% 06-diseno-arquitectura.tex — Fuente: docs/informe-final/06-diseno-arquitectura.md

\chapter{Diseño del sistema y arquitectura}
\label{cap:diseno-arquitectura}

\section{Modelo C4}

\subsection{Nivel 1 --- Contexto}

Artisync se sitúa como sistema central que conecta tres actores (Cliente,
Creador/Artista, Administrador) con tres sistemas externos: PayPal
(pagos, patrón \emph{escrow}), un servicio de IA de verificación de
identidad/certificados, y un proveedor de almacenamiento de objetos
(Azure Blob Storage o local, según \texttt{ADR-007}). El modelo C4 usado
como base sigue a Brown~\cite{BROWN2018}. Diagrama completo:
\texttt{docs/diagramas/C4\_Nivel1\_Contexto.md} y
\texttt{C4-nivel1-context\_diagram.svg}.

\subsection{Nivel 2 --- Contenedores}

Cuatro contenedores desplegables vía Docker Compose: \textbf{frontend}
(Angular, SPA servida por nginx), \textbf{backend} (Spring Boot, API REST
documentada con OpenAPI/Springdoc), \textbf{postgres} (PostgreSQL~16,
esquema versionado con Flyway) y \textbf{redis} (caché de catálogo,
blacklist de JWT revocados, rate limiting, tickets 2FA). Diagrama
completo: \texttt{docs/diagramas/C4\_Nivel2\_Contenedores.md}.

Como se ilustra en la Figura~\\ref{fig:c4-nivel2}, 
Como se ilustra en la Figura~\\ref{fig:mer}, 
\begin{figure}[H]
\centering
\includegraphics[width=0.85\textwidth]{c4-nivel2-contenedores.png}
\caption{Diagrama C4 Nivel 2 --- Contenedores de Artisync. Fuente:
\texttt{docs/diagramas/c4-nivel2-contenedores.png}.}
\label{fig:c4-nivel2}
\end{figure}

\textbf{Nota de honestidad sobre versiones documentadas vs. reales:} los
diagramas C4 (y \texttt{ADR-001}) documentan la pila como ``Java~25 +
Spring Boot~4.0.1'' y ``PostgreSQL~18/Angular~22'' en su versión
original. La pila \textbf{realmente construida y desplegada}, verificada
contra \texttt{pom.xml}, los \texttt{Dockerfile} y
\texttt{docker-compose.yml} a la fecha de este documento, es: \textbf{Java
21~LTS} (target de bytecode) \textbf{+ Spring Boot~4.1.0}, compilado en
contenedor \texttt{maven:3.9-eclipse-temurin-21} y ejecutado en
\texttt{eclipse-temurin:21-jre-alpine}; \textbf{Angular \^{}22.0.0} sobre
\texttt{node:24-alpine}; \textbf{PostgreSQL~16} (\texttt{postgres:16},
anclado por digest \texttt{sha256}) y \textbf{Redis~7}
(\texttt{redis:7-alpine}, igualmente anclado). La documentación
arquitectónica no se actualizó en cascada cuando el equipo fijó las
versiones finales --- ver \texttt{docs/entorno/versions.txt} para la
versión verificada y autoritativa, que este capítulo usa como fuente en
vez de repetir la cifra desactualizada de los diagramas.

\subsection{Nivel 3 --- Componentes (backend)}

El backend se organiza en siete módulos funcionales por dominio
(\texttt{seguridad}, \texttt{perfil}, \texttt{catalogo},
\texttt{pedido}/flujo de trabajo, \texttt{legal}, \texttt{comunicacion},
\texttt{social}), cada uno con sus capas
\texttt{entity}/\texttt{dto}/\texttt{repository}/\texttt{service}/\texttt{controller}.
Detalle completo: \texttt{docs/diagramas/C4\_Nivel3\_Componentes\_Backend.md}.

\section{Resumen de los siete registros de decisión de arquitectura (ADR)}

Los siete ADR siguen la plantilla de Nygard~\cite{NYGARD2011}:

Como se detalla en la Tabla~\\ref{tab:adrs}, 
Como se detalla en la Tabla~\\ref{tab:iso25010}, 
\begin{table}[H]
\centering
\caption{Resumen de los siete ADR de Artisync}
\label{tab:adrs}
\begin{tabular}{@{}p{1.4cm}p{3.4cm}p{9.4cm}@{}}
\toprule
\textbf{ADR} & \textbf{Título} & \textbf{Decisión resumida} \\
\midrule
ADR-001 & Selección de la pila tecnológica principal & Java/Spring Boot en
  el backend (mayor throughput, seguridad/persistencia integradas),
  Angular/TypeScript en el frontend, PostgreSQL como motor relacional,
  Redis como caché, Docker Compose como orquestador --- ver nota de
  versiones \S6.1. \\
ADR-002 & Esquema de autenticación y control de acceso & JWT (HS256)
  devuelto en el cuerpo de la respuesta para el access token (no en
  cookie); refresh token y ticket pre-auth de 2FA en cookies
  \texttt{HttpOnly + SameSite=Strict}; blacklist de JTI en Redis; bcrypt
  factor 12; RBAC vía Spring Security. Revisado en agosto 2026 tras
  corregir una afirmación incorrecta de una versión anterior del ADR (ver
  texto completo). \\
ADR-003 & Selección del gestor de base de datos & PostgreSQL~16 sobre
  MySQL/MariaDB y MongoDB, por ACID completo, tipos \texttt{DECIMAL}
  precisos para montos financieros, y soporte maduro de PL/pgSQL para la
  estrategia híbrida de acceso a datos (ADR-006). \\
ADR-004 & Estrategia de caché & Redis~7 para caché de catálogo (TTL
  configurable), blacklist de JTI, cuotas de rate limiting y tickets
  pre-auth 2FA, sobre caché en memoria de la JVM, por ser compartido entre
  réplicas si el backend escala horizontalmente. \\
ADR-005 & Estrategia de despliegue y reproducibilidad & Docker Compose con
  imágenes ancladas por digest \texttt{sha256} + \texttt{Makefile} de un
  solo comando, sobre despliegue manual o imágenes por tag variable, para
  máxima reproducibilidad (badge ACM \emph{Artifacts Evaluated ---
  Reusable}). Las tres acciones pendientes de la Tercera Entrega
  (Frontend/ ausente, digests por tag, Makefile inexistente) están
  resueltas y verificadas a la fecha de la Entrega Final. \\
ADR-006 & Estrategia híbrida de acceso a datos & CRUD elementales vía
  JPA/Spring Data; operaciones con joins, agregaciones, reportes,
  actualizaciones masivas o validaciones cruzadas vía
  procedimientos/funciones PL/pgSQL versionados en \texttt{db/procs/},
  invocados exclusivamente mediante JPA~2.1. Ampliado el 16 de agosto de
  2026 con 7 rutinas nuevas conectadas end-to-end (ver
  Capítulo~\ref{cap:implementacion}). \\
ADR-007 & Almacenamiento de archivos en la nube & Azure Blob Storage como
  único destino en la nube (sobre un esquema híbrido Cloudinary+Azure Blob
  o mantener el volumen local), detrás de la interfaz
  \texttt{AlmacenamientoDocumentos}, seleccionable por configuración
  (\texttt{documentos.proveedor}). \\
\bottomrule
\end{tabular}
\end{table}

Texto completo de cada ADR, con contexto, opciones descartadas y
consecuencias, en \texttt{docs/adr/}.

\section{Atributos de calidad (ISO/IEC 25010) --- prioridad, escenario y estrategia}

\begin{table}[H]
\centering
\caption{Atributos de calidad ISO/IEC~25010~\cite{ISO25010} priorizados}
\label{tab:iso25010}
\begin{tabular}{@{}p{3.0cm}p{1.6cm}p{4.0cm}p{5.6cm}@{}}
\toprule
\textbf{Característica ISO~25010} & \textbf{Prioridad} & \textbf{Escenario} & \textbf{Estrategia aplicada} \\
\midrule
Eficiencia de desempeño (rendimiento) & Alta & Listado de catálogo bajo 50
  usuarios concurrentes & Caché Redis con TTL (ADR-004); índices en
  columnas de filtro; medido con k6 (Capítulo~\ref{cap:evaluacion-resultados}) \\
Seguridad & Alta & Acceso no autorizado a rutas protegidas, inyección
  SQL, credenciales en tránsito & RBAC + JWT (ADR-002);
  procedimientos/consultas parametrizadas (ADR-006); rate limiting;
  auditoría estática y dinámica automatizada
  (Capítulo~\ref{cap:evaluacion-resultados}) \\
Fiabilidad & Alta & Caída de un contenedor dependiente (Redis, Postgres) &
  \texttt{healthcheck} de Docker Compose con
  \texttt{depends\_on: condition: service\_healthy};
  \texttt{/actuator/health} expone estado por componente \\
Usabilidad & Media & Usuario nuevo completa el flujo de contratación sin
  asistencia & Formularios dinámicos, flujo de contratación acotado (SRS
  \texttt{REQ-NF-008}); medido con SUS
  (Capítulo~\ref{cap:evaluacion-resultados}) \\
Mantenibilidad & Media & Incorporar un nuevo procedimiento almacenado sin
  tocar el código Java existente & Módulos desacoplados por dominio
  (Nivel~3 C4); \texttt{scripts/sync-procs.sh} mantiene sincronizada la
  migración repetible \\
Portabilidad & Media & Desplegar en un proveedor cloud distinto al usado
  en desarrollo & Configuración externalizada (\texttt{.env}),
  \texttt{documentos.proveedor} seleccionable (ADR-007), imágenes ancladas
  por digest (ADR-005) \\
Compatibilidad & Baja & Integración con PayPal Orders v2 y servicio
  externo de IA & Contratos de API versionados, adaptadores dedicados por
  integración \\
\bottomrule
\end{tabular}
\end{table}

\section{Modelo de datos}

El esquema relacional completo se versiona con Flyway; el modelo
entidad-relación se mantiene en \texttt{docs/diagramas/Entidad\_Relacion.png}
(ver nota de conformidad de nombres de tablas en \texttt{ADR-003}).

\begin{figure}[H]
\centering
\includegraphics[width=0.95\textwidth]{Entidad_Relacion.png}
\caption{Modelo entidad-relación de Artisync. Fuente:
\texttt{docs/diagramas/Entidad\_Relacion.png}.}
\label{fig:mer}
\end{figure}

\section{Matriz de trazabilidad}

La matriz completa (37 filas: requisito $\to$ historia de usuario $\to$
caso de uso $\to$ módulo de código $\to$ endpoint $\to$ prueba
automatizada $\to$ \textbf{tipo de acceso} $\to$ evidencia empírica
$\to$ estado) está en \texttt{docs/trazabilidad/matriz.csv}, validada
automáticamente en CI por \texttt{scripts/validate-traceability.sh}. La
columna \texttt{tipo\_acceso} distingue \texttt{CRUD-ORM} de \texttt{SP}
(procedimiento almacenado): a la fecha de este documento, 8 filas están
marcadas \texttt{SP} (correspondientes a los 7 requisitos que ya invocan
una de las rutinas conectadas end-to-end, ver
Capítulo~\ref{cap:implementacion}), el resto \texttt{CRUD-ORM}.
